\section{Limitations}

Several limitations define the interpretation of the study. First, the design is observational. The analyses do not identify causal effects of CMAT use, instructor movement, academic failure, or degree programme. Students self-select into formal support and, in later contexts, into course sections subject to constraints that are only partly observed.

Second, the first-MU experience states use realized academic outcomes. Classroom-relative performance and classroom pass-rate context are measured at the end of the same course period in which CMAT visits occur. Their relationship with CMAT attendance is therefore contemporaneous. The analysis can show that certain combinations of realized individual and classroom outcomes coexist with higher or lower formal support use, but it cannot establish whether difficulty caused attendance, attendance altered performance, or both were shaped by earlier unobserved factors.

Third, the later Calculus instructor set is only partially observed. We know which instructors taught the course in each period and which have sufficient historical data to be ranked, but we do not observe each student's feasible timetable, section capacity, registration constraints, or the information available to the student at enrolment. Consequently, later instructor context is an observed enrolment outcome, not a revealed-preference measure.

Fourth, historical instructor outcomes are not intrinsic instructor characteristics. Prior pass rates and mean grades may reflect student composition, curricular changes, assessment design, course timing, and other contextual factors. Strictly prior measurement avoids using future information to characterize a focal enrolment, but it does not turn historical outcomes into causal measures of difficulty or quality.

Fifth, CMAT records capture only one form of formal help-seeking. Students may seek assistance from classmates, friends, instructors, online resources, private tutors, or family. The broader help-seeking literature shows that sources and goals matter \citep{fong2023}, while first-year research demonstrates that social and academic support are embedded in wider networks \citep{wilcox2005}. Non-use of CMAT therefore cannot be classified as disengagement.

Finally, programme comparisons remain descriptive. Degree programme may proxy curricular intensity, student selection, prior preparation, peer norms, or programme-specific advising. The omnibus tests establish heterogeneity but do not isolate a disciplinary causal mechanism.

\draftnote{Add the institutionally approved ethics/data-governance statement before submission. Do not infer or invent an approval number.}

\section{Conclusion}

This study uses longitudinal administrative records to examine academic challenge, formal help-seeking, and adaptation without treating those constructs as interchangeable. The central empirical pattern is a mismatch between academic need and formal support use. Students who face difficult classroom environments while maintaining relatively stronger performance are the most likely to use the mathematics support centre, whereas students with substantially weaker relative outcomes use formal support much less. The pattern is stable across alternative definitions of academic strain.

Later adaptation is similarly heterogeneous. A difficult MU experience does not generally lead students to enrol with historically higher-outcome Calculus instructors. Immediate failure, however, is followed by near-universal instructor change and, for most comparable students, movement toward an instructor associated with stronger prior academic outcomes. Formal support use does not increase in parallel. Degree programme adds further, but modest, heterogeneity to these behavioural traces.

The findings support a restrained use of administrative data in engagement research. Attendance records, grades, course repetition, and instructor enrolment can reveal meaningful sequences of behaviour, but none of them alone measures student engagement. Their value lies in preserving their differences and studying how they unfold over time. In this setting, the resulting sequence is not a simple path from difficulty to help-seeking to success. It is a set of partially independent responses to academic challenge, shaped by institutional context and disciplinary environment.


\bibliographystyle{unsrtnat}
\bibliography{references}

\clearpage
\appendix
\section{Supporting descriptive tables}
\label{app:supporting-tables}
\setcounter{table}{0}
\renewcommand{\thetable}{A\arabic{table}}

\begin{landscape}
\begin{table}[p]
\centering
\small
\caption{Primary first-MU experience states and subsequent outcomes}
\label{tab:states}
\begin{tabular}{@{}lrrrrr@{}}
\toprule
State & $N$ & MU CMAT & Eventual MU pass & Later Calculus & Calc CMAT \\
\midrule
Lower strain & 3,679 & 18.8\% & 99.8\% & 77.5\% & 23.6\% \\
Contextual challenge & 1,012 & 28.6\% & 97.7\% & 76.3\% & 30.6\% \\
Individual strain & 673 & 9.5\% & 58.5\% & 42.2\% & 13.7\% \\
Compounded strain & 162 & 8.0\% & 30.9\% & 16.0\% & 15.4\% \\
Adverse/non-numeric & 1,101 & 16.2\% & 44.2\% & 19.7\% & 15.2\% \\
\bottomrule
\end{tabular}
\caption*{\footnotesize Notes: MU CMAT is any recorded support-centre use in the first MU period. Later Calculus is the proportion observed in a later Calculus attempt under the study's coverage rules. Calc CMAT is support use among those observed in later Calculus. Experience states are defined independently of CMAT attendance.}
\end{table}
\end{landscape}

\begin{landscape}
\begin{table}[p]
\centering
\small
\caption{Observed adaptation after failed MU attempts}
\label{tab:failure}
\begin{tabular}{@{}lrrrrrr@{}}
\toprule
Failed attempt & $N$ & Change prof. & Prior CMAT & Next CMAT & Next pass & Move higher \\
\midrule
1 & 820 & 97.0\% & 18.9\% & 9.8\% & 64.6\% & 62.3\% \\
2 & 143 & 93.0\% & 3.5\% & 4.9\% & 54.5\% & 60.0\% \\
3 & 30 & 100.0\% & 6.7\% & 3.3\% & 40.0\% & 63.6\% \\
\bottomrule
\end{tabular}
\caption*{\footnotesize Notes: ``Move higher'' is calculated among transitions with comparable strictly-prior instructor-outcome percentiles ($N=387$, 85, and 22 for failed attempts 1--3). Later-attempt samples are small and are reported descriptively.}
\end{table}
\end{landscape}

\begin{landscape}
\begin{table}[p]
\centering
\small
\caption{Incremental contribution of degree programme}
\label{tab:degree}
\begin{tabular}{@{}lrrrr@{}}
\toprule
Outcome & $N$ & Programme levels & $\Delta R^2$ & Joint $p$ \\
\midrule
First-MU CMAT use & 6,627 & 30 & 0.0154 & $<.001$ \\
Later Calculus CMAT use & 4,151 & 25 & 0.0250 & $<.001$ \\
Later instructor historical rank & 3,224 & 25 & 0.0143 & $<.001$ \\
\bottomrule
\end{tabular}
\end{table}
\end{landscape}

\draftnote{Prepare supplementary material mapping the manuscript claims to reviewed tables 220--233 and include the full threshold, choice-set, and programme-interaction sensitivities.}

\end{document}
